\documentclass[10pt,a4paper]{amsart}
\usepackage[margin=19mm]{geometry}
\usepackage{amsmath,amssymb,mathtools}
\usepackage[scr=rsfs,cal=euler]{mathalfa}
\usepackage{xfrac}
\usepackage[unicode,hidelinks,bookmarks=false]{hyperref}
\newcommand{\ie}{i.e.\ }
\newcommand{\cf}{cf.\ }
\newcommand{\resp}{respectively\ }
\newcommand{\Subsection}{\subsection}
\providecommand{\nobreakdash}{\nobreak\hbox{-}}
\setlength{\emergencystretch}{2em}
\multlinegap0pt
\usepackage{bm}
\usepackage{bbm}
\usepackage{dsfont}
\usepackage{xspace}
\usepackage{cancel}
\usepackage{mathtools}
\usepackage{amsthm}
\makeatletter
\@ifundefined{Theorem}{\newtheorem{Theorem}{Theorem}}{}
\@ifundefined{Lemma}{\newtheorem{Lemma}[Theorem]{Lemma}}{}
\makeatother
\title[Factorization method for a clamped obstacle from near-field ]{Factorization method for a clamped obstacle from near-field measurements via a far-field transformation}
\author{General Ozochiawaeze, Isaac Harris}
\date{Source: arXiv:2609.20528v1 (CC BY 4.0)}
\begin{document}
\maketitle

\noindent\emph{Source and licence.} Factorization method for a clamped obstacle from near-field measurements via a far-field transformation. Version arXiv:2609.20528v1 (CC BY 4.0), distributed under the Creative Commons Attribution 4.0 International licence, as recorded in the arXiv licence metadata for this exact version and shown on the abstract page https://arxiv.org/abs/2609.20528v1. Full rights notice: licenses/arxiv-2609.20528-CC-BY-4.0.txt.\par\smallskip

\noindent\textit{Editorial scope.} Counted unit: the Lemma whose source label is \texttt{None} in the article (source0.tex lines 397--402, source label \texttt{None}), delivered together with the article's own complete original proof of it (source0.tex lines 403--476, one \texttt{proof} environment of 74 lines). No mathematics is written, repaired or re-derived: statement and proof are the article's own text line for line. Required context retained uncounted: eqnbcs (lines 160--162); SRCs (lines 168--170); NF\_operator (lines 210--228); ext\_BVP (lines 269--272); clamped\_bcs\_inhomog (lines 274--277). Every definition, display and label of those lines is retained unchanged. Omitted from this package and not redistributed: 1--159, 163--167, 171--209, 229--268, 273--273, 278--396, 477--2342. Nothing inside a retained excerpt is omitted or altered: every retained body is delivered exactly as the article's lines are written, so no omission receipt of any kind accompanies this package. Citations: the retained excerpts carry no citation command at all, so no bibliography entry of the article is transcribed and no reference is redistributed. Reference adaptation: none was needed and none was made; every reference inside the delivered unit resolves inside it, so no unresolved reference or citation is produced. Printed slips kept literal: no printed word, symbol, hypothesis or number of the counted statement or of its proof was altered. Layout and class: the article's own class is \texttt{article} with packages including amsmath, amsfonts, amssymb, amsthm, color, graphicx, mathtools, cite, enumitem, lipsum, algorithm, algorithmic, dsfont, hyperref; the wrapper is the lane's pinned \texttt{amsart} editorial wrapper at 10pt on A4, which restates the theorem-like environments the retained text needs and adds the article's own macro definitions that the retained text needs (none), with extra wrapper packages (bm, bbm, dsfont, xspace, cancel, mathtools). The delivered numbering is the unit's own. Assets: no retained span contains a figure environment, a graphics-inclusion command, a PDF-page inclusion, a table or a TikZ picture; no image file is copied into this package, the package declares no asset dependency and no image file is redistributed with it. Licence: version arXiv:2609.20528v1 of this article is distributed under the Creative Commons Attribution 4.0 International licence, as recorded in the arXiv licence metadata for this exact version and shown on the abstract page https://arxiv.org/abs/2609.20528v1. A full rights notice is delivered at licenses/arxiv-2609.20528-CC-BY-4.0.txt. Characters: every retained excerpt is pure ASCII, so the delivered engine, XeLaTeX with its default Latin Modern text font, reports no missing character.
\begin{align}\label{eqnbcs}
    \Delta^2u^s-\kappa^4u^s=0,\quad x\in\mathbb R^2\setminus\overline D
\end{align}

\begin{align}\label{SRCs}
    \lim_{r\to\infty} \sqrt{r}(\partial_r u^s-\mathrm{i}\kappa u^s)=0\quad \text{and}\quad \lim_{r\to\infty} \sqrt{r}(\partial_r \Delta u^s-\mathrm{i}\kappa \Delta u^s)=0, \quad r=|x|
\end{align}

\begin{equation}\label{NF_operator}
\mathcal N
\begin{pmatrix}
h\\
t
\end{pmatrix}(x)
=
\begin{pmatrix}
\displaystyle
\int_{\mathcal C}
\big(u^{s}(x,y)\,h(y) + \tilde u^{s}(x,y)\,t(y)\big)\,ds(y)
\\[2ex]
\displaystyle
\partial_{\nu_x}
\int_{\mathcal C}
\big(u^{s}(x,y)\,h(y) + \tilde u^{s}(x,y)\,t(y)\big)\,ds(y)
\end{pmatrix},
\qquad x \in \mathcal C.
\end{equation}

\begin{equation}\label{ext_BVP}
\Delta^2 v - \kappa^4 v = 0 
\qquad \text{in } \mathbb{R}^2 \setminus \overline{D},
\end{equation}

\begin{align}\label{clamped_bcs_inhomog}
v = \phi \quad\text{and}\quad \partial_{\nu} v = \psi 
\qquad \text{on } \partial D.
\end{align}

\begin{Lemma}
The near-field operator $\mathcal N$ admits the factorization
\[
\mathcal N=-\mathcal U\mathcal B^\top .
\]
\end{Lemma}

\begin{proof}
Let $(h,t)^\top\in [L^2(\mathcal C)]^2$ and define the incident field
\[
v^i(x)=
\int_{\mathcal C}
\big(
\mathbb G(x,y)h(y)
+\partial_{\nu_y}\mathbb G(x,y)t(y)
\big)\,ds(y).
\]
By the definition of $\mathcal B^\top$, the Cauchy data of $v^i$ on
$\partial D$ satisfy
\[
\begin{pmatrix}
v^i\\
\partial_\nu v^i
\end{pmatrix}
=
\mathcal B^\top
\begin{pmatrix}
h\\
t
\end{pmatrix}.
\]
Let $v$ denote the unique radiating solution to \eqref{ext_BVP}-\eqref{clamped_bcs_inhomog} with clamped boundary conditions given by 
\[
v=-v^i,\qquad
\partial_\nu v=-\partial_\nu v^i
\quad\text{on }\partial D .
\]
Therefore, by the definition of $\mathcal U$,
\[
\mathcal U\mathcal B^\top
\begin{pmatrix}
h\\
t
\end{pmatrix}
=
-
\begin{pmatrix}
v^s|_{\mathcal C}\\
\partial_\nu v^s|_{\mathcal C}
\end{pmatrix}.
\]
On the other hand, {by well-posedness of the forward problem \eqref{eqnbcs}-\eqref{SRCs} and superposition}, the scattered field $v$ admits the representation
\[
v(x)
=
\int_{\mathcal C}
\big(
u^s(x,y)h(y)+\tilde u^s(x,y)t(y)
\big)\,ds(y),
\]
where $u^s(\cdot,y)$ and $\tilde u^s(\cdot,y)$ are the scattered fields
generated by the point-source incident fields
$\mathbb G(\cdot,y)$ and $\partial_{\nu_y}\mathbb G(\cdot,y)$,
respectively. Hence, using the definition of $\mathcal N$ in
\eqref{NF_operator}, we obtain
\[
\mathcal U\mathcal B^\top
\begin{pmatrix}
h\\
t
\end{pmatrix}
=
-\mathcal N
\begin{pmatrix}
h\\
t
\end{pmatrix}
\quad \text{ which gives that } \quad \mathcal N=-\mathcal U\mathcal B^\top 
\]
proving the claim. 
\end{proof}

\end{document}
